% Generated semantic source for AI Almanach v3.
% Edit v3 directly after initial generation; do not overwrite
% editorial changes by rerunning the converter casually.

% ==== Start-up ================================================================
\chapter{Building an AI Venture}
\chapterlead{How can an AI venture test problem, feasibility, and economics?}{Problem}{Experiment}{Economics}{Scale}{startup}

\label{ch:startup}

\section{From an Interesting Model to a Valuable Product}

\begin{v3topic}{Start with a Painful Workflow}
A start-up begins with a costly, frequent, and observable problem---not with a
model looking for somewhere to live.
Describe who performs the work, what triggers it, which inputs are available,
what outcome matters, and what happens when the decision is wrong.
Customer interviews are evidence about behaviour only when they reconstruct
specific past events; compliments and hypothetical purchase intent are weak
signals\textsuperscript{\cite{constableTalkingHumansSuccess2014}}.

\end{v3topic}
\begin{v3topic}{Write the Non-AI Baseline First}
Before training, compare the proposed system with a rule, search interface,
spreadsheet, process redesign, or human service.
The baseline exposes whether uncertainty-tolerant statistical behaviour is
actually needed.
It also establishes the minimum quality, cost, latency, and reliability that
an AI solution must beat.

\end{v3topic}
\begin{v3topic}{Separate Four Hypotheses}
An AI product simultaneously tests a \keyterm{problem hypothesis} (the pain is
real), a \keyterm{value hypothesis} (improvement changes behaviour), a
\keyterm{feasibility hypothesis} (data and technology meet the target), and a
\keyterm{viability hypothesis} (the economics and obligations are sustainable).
A successful model experiment validates only part of the feasibility
hypothesis.
It does not prove demand, workflow adoption, lawful data rights, or a viable
business.
\end{v3topic}

\begin{cautionbox}[The benchmark-to-business fallacy]
An improvement on a public benchmark may have no economic value if the error
being reduced is cheap, the workflow cannot use the output, integration costs
dominate, or the data distribution differs from the customer environment.
Define the operational decision and its error costs before choosing a metric.
\end{cautionbox}

\section{Lean Startup: The Method and Its Evidence}
\label{sec:startup-lean}

\subsection{The Reframe}

Eric Ries's \emph{The Lean
Startup}\textsuperscript{\cite{riesLeanStartup2011}} --- published in German
as \emph{Lean Startup: Schnell, risikolos und erfolgreich Unternehmen
gründen}\textsuperscript{\cite{riesLeanStartupDeutsch2014}} --- is the single
most influential text on how to build a company under uncertainty, and its
central move is a definition:

\begin{quote}
\itshape A startup is a human institution designed to create a new product or
service under conditions of extreme uncertainty.
\end{quote}

Three consequences follow, and they are the whole method. If the defining
condition is uncertainty, then the scarce resource is not capital or code but
\keyterm{validated learning} --- knowledge about customers demonstrated by
evidence rather than assumed. If learning is the unit of progress, then
progress can be \emph{managed}, and entrepreneurship becomes a management
discipline rather than a temperament. And if learning is what you are buying,
then any effort that does not produce it is \keyterm{waste}, however
impressive the artefact it produces.

\begin{plainlanguage}
The ordinary way to start a company is to decide what to build, build it, and
find out at launch whether anyone wanted it. Lean startup inverts that: work
out what you are \emph{assuming}, find the cheapest experiment that could show
the assumption is wrong, and run it before building. A startup is then not a
small version of a big company --- it is a machine for converting money into
answers, and the question is how many answers per euro it produces.
\end{plainlanguage}

\begin{engineernote}
The method is close kin to everything this book says about evaluation. An
assumption is a hypothesis; an MVP is an experiment; a cohort is a held-out
sample; a vanity metric is training accuracy on data you already fit. If you
have internalised \cref{ch:evaluation-mlops}, you already have the reflexes ---
lean startup applies them to the business rather than to the model.
\end{engineernote}

\subsection{Vision: The Two Hypotheses and Validated Learning}

Ries divides the book into three parts --- \emph{Vision}, \emph{Steer}, and
\emph{Accelerate} --- and Part~I asks what is actually being tested. Every new
venture rests on two leaps of faith:

\begin{description}
  \item[The value hypothesis] Does the product deliver value to customers once
    they use it? Not ``would you use this?'' but observed behaviour.
  \item[The growth hypothesis] How will new customers discover the product,
    and does that mechanism compound?
\end{description}

These are separable and must be tested separately, in that order. A product
that grows without delivering value is a leaking bucket filled faster; a
product that delivers value without a growth mechanism is a good business that
stays small.

\subsection{Steer: Build--Measure--Learn}

The operational core is a feedback loop --- and the most-quoted, least-applied
insight about it is that you \emph{plan} it backwards while you \emph{run} it
forwards.

\begin{figure}[htbp]
\centering
\begin{tikzpicture}[
  node distance=13mm,
  bx/.style={draw=AlmanachTeal,fill=AlmanachMist,rounded corners=1mm,
    minimum height=9mm,minimum width=24mm,font=\sffamily\scriptsize,
    align=center},
  art/.style={font=\sffamily\scriptsize,color=AlmanachOchre!80!black},
  fw/.style={-{Stealth[length=2mm]},thick,draw=AlmanachTeal},
  bw/.style={-{Stealth[length=2mm]},thick,draw=AlmanachRed,dashed},
  lbl/.style={font=\sffamily\scriptsize,color=AlmanachInk!70}]
  \node[bx] (ideas) at (0,16mm) {IDEAS};
  \node[bx] (build) at (34mm,0) {BUILD};
  \node[bx] (code) at (34mm,-14mm) {\textcolor{AlmanachOchre!80!black}{product}};
  \node[bx] (measure) at (0,-28mm) {MEASURE};
  \node[bx] (data) at (-34mm,-14mm) {\textcolor{AlmanachOchre!80!black}{data}};
  \node[bx] (learn) at (-34mm,0) {LEARN};
  \draw[fw] (ideas) -- (build);
  \draw[fw] (build) -- (code);
  \draw[fw] (code) -- (measure);
  \draw[fw] (measure) -- (data);
  \draw[fw] (data) -- (learn);
  \draw[fw] (learn) -- (ideas);
  \draw[bw] (-16mm,20mm) .. controls (20mm,26mm) and (46mm,10mm) .. (46mm,-6mm);
  \node[lbl,color=AlmanachRed,anchor=west] at (-14mm,25mm)
    {plan in this direction};
  \node[lbl,color=AlmanachTeal,anchor=north,align=center] at (0,-36mm)
    {execute in this direction};
  \node[lbl,anchor=west,align=left,text width=40mm] at (52mm,-14mm)
    {\emph{Minimise total time through the loop} --- that is the only thing
     the loop optimises};
\end{tikzpicture}
\caption{The build--measure--learn loop. Design it backwards: decide what you
need to \emph{learn}, then what you must \emph{measure} to learn it, then the
smallest thing you can \emph{build} to produce that measurement. Teams that
plan forwards build a product and then look for something to measure about
it.}
\label{fig:startup-bml-loop}
\end{figure}

\begin{cautionbox}[What a minimum viable product is, and what it is not]
The \keyterm{MVP} is the version of the product that enables a full turn of
the loop with the least effort. It is an \emph{experimental instrument}, and
several common readings are wrong:
\begin{itemize}
  \item It is \textbf{not} version 1.0 of the real product, nor a
    stripped-down but complete thing you would be proud to ship.
  \item It is \textbf{not} necessarily software. A landing page, a manual
    service, or a slide deck can be an MVP if it produces the measurement.
  \item It \textbf{is} allowed to be embarrassing --- Ries's point is that the
    embarrassment is the price of learning early.
  \item It \textbf{must} be tied to a metric decided in advance. An MVP with
    no pre-declared measurement is just a small product.
\end{itemize}
Useful forms: the \keyterm{concierge MVP} (deliver the service by hand to a
handful of customers), the \keyterm{Wizard-of-Oz MVP} (the interface looks
automated; humans do the work behind it), the \keyterm{smoke test} (sell it
before it exists and count who tries to buy), and the single-feature product.
\end{cautionbox}

\subsection{Innovation Accounting}

If learning is progress, it needs a ledger. \keyterm{Innovation accounting}
replaces gross totals with three milestones: establish a baseline with a first
MVP; tune the engine by moving the metric from baseline toward the ideal; then
decide to pivot or persevere. Metrics must be \emph{actionable} (show clear
cause and effect), \emph{accessible} (readable by the people who must act),
and \emph{auditable} (traceable to real customer behaviour).

\begin{workedexample}[A vanity metric hiding a deteriorating product]
A dashboard shows total registered users growing from $5{,}000$ to $40{,}000$
in six months. The board is delighted. Now split the same data by
\keyterm{cohort} --- grouping users by the month they joined and following each
group forward:

\smallskip
\begin{tblr}{
  width=0.92\linewidth,
  colspec={Q[l,0.16\linewidth] Q[c] Q[c] Q[c]},
  row{1}={font=\bfseries,bg=AlmanachMist},
  hlines,vlines}
Cohort & Signups & Converted to paid & Rate \\
Month 1 & 1{,}000 & 42 & $4.2\,\%$ \\
Month 2 & 1{,}800 & 70 & $3.9\,\%$ \\
Month 4 & 4{,}500 & 158 & $3.5\,\%$ \\
Month 6 & 8{,}000 & 248 & $3.1\,\%$ \\
\end{tblr}
\smallskip

\emph{Read both readings.} Absolute paying customers rose from $42$ to $248$ ---
a genuine $5.9\times$. Conversion \emph{per cohort} fell from $4.2\,\%$ to
$3.1\,\%$, a relative decline of $26\,\%$.

\emph{What is actually happening.} Growth is being bought, and each new
customer is worth less than the last. Every feature shipped in those six
months either failed to improve conversion or actively hurt it, and the total
line concealed that completely. Had the team been steering by the total, they
would have concluded the product was working.

\emph{The rule.} Cumulative totals only go up and therefore cannot falsify
anything. Steer by cohort behaviour and by split tests, which can.
\end{workedexample}

\subsection{Pivot or Persevere}

A \keyterm{pivot} is a structured course correction that tests a new
fundamental hypothesis --- not a synonym for a change of plan, and not a
failure. Ries catalogues ten forms.

\begin{table}[htbp]
\centering
\caption{The ten pivot types. The value of the taxonomy is diagnostic: naming
which hypothesis failed tells you which pivot is available, and stops a team
changing everything at once.}
\label{tab:startup-pivot-types}
\begin{tblr}{
  width=\textwidth,
  colspec={Q[l,0.20\textwidth] X[l]},
  row{1}={font=\bfseries,bg=AlmanachMist},
  hlines,vlines}
Pivot & What changes \\
Zoom-in & One feature becomes the whole product \\
Zoom-out & The whole product becomes one feature of something larger \\
Customer segment & The product roughly works, but for a different buyer \\
Customer need & Same buyer, different problem worth solving \\
Platform & From application to platform, or the reverse \\
Business architecture & High margin/low volume $\leftrightarrow$ low margin/high volume \\
Value capture & A different revenue or pricing model \\
Engine of growth & Switch between viral, sticky, and paid growth \\
Channel & A different route to the customer \\
Technology & Same problem and customer, different technical means \\
\end{tblr}
\end{table}

\begin{engineernote}
Schedule the pivot-or-persevere decision as a recurring meeting with a fixed
cadence, and set the criteria before you need them. Ries's observation is that
most startups fail to pivot not because the evidence is unclear but because
nobody is obliged to look at it --- and because a vanity metric always offers a
reason to persevere. This is the same discipline as the kill criteria later in
this chapter, and the same discipline as a pre-registered evaluation
threshold in \cref{ch:evaluation-mlops}.
\end{engineernote}

\subsection{Accelerate: Engines of Growth and Small Batches}

Part~III asks how to run the loop faster and how growth compounds. Ries
identifies three \keyterm{engines of growth}, each with its own governing
metric, and warns against pursuing more than one at a time.

\begin{workedexample}[The engines, and the arithmetic that decides them]
\emph{Viral engine.} Growth comes from use itself. The governing quantity is
the viral coefficient $k = i \times c$, where $i$ is invitations sent per user
and $c$ the fraction that convert.
\[
  i = 2.5,\ c = 0.15 \;\Rightarrow\; k = 0.375,
  \qquad
  \text{total users per initial user } = \frac{1}{1-k} = 1.6 .
\]
The campaign dies out. Now
\[
  i = 4,\ c = 0.30 \;\Rightarrow\; k = 1.2 > 1 \;\Rightarrow\;
  \text{unbounded growth.}
\]
Note the cliff at $k=1$: at $k = 0.9$ each initial user yields $10$ in total
--- large, but finite and self-extinguishing. Between $0.9$ and $1.2$ lies the
difference between a campaign and a company.

\emph{Sticky engine.} Growth comes from retention. The governing quantity is
the compounding rate $=$ acquisition rate $-$ churn rate. At $6\,\%$
acquisition and $4\,\%$ churn per month:
\[
  1.02^{12} = 1.27 \;\Rightarrow\; 27\,\% \text{ annual growth.}
\]
At $6\,\%$ acquisition and $8\,\%$ churn the business shrinks, no matter how
much is spent on marketing --- churn must be fixed first.

\emph{Paid engine.} Growth is bought from margin. It works only while
lifetime value exceeds acquisition cost with enough headroom to fund the next
cohort; the unit economics of \cref{sec:startup-unit-economics} are exactly
this calculation.

\emph{Why one at a time.} Each engine has a different governing metric and
implies different work. A team optimising all three optimises none, and cannot
attribute any change it observes.
\end{workedexample}

Two further accelerators complete Part~III. \keyterm{Small batches} shorten
the loop: work in the smallest increment that can be measured, because a large
batch hides which change caused which effect --- the same argument as changing
one variable at a time in an experiment. And the \keyterm{Five Whys} converts
each failure into a proportional investment: ask ``why'' five times to move
from a symptom to a root cause, then fix at every level in proportion to the
damage, which prevents both over-reaction and repeated incidents.

% ===========================================================================
\section{Does Lean Startup Work? The Evidence}
% ===========================================================================

Lean startup is a management theory, and this book's standard is that
management theories, like models, are claims to be tested. The evidence is
better than most business methods can offer and weaker than its popularity
implies.

\begin{table}[htbp]
\centering
\caption{Empirical findings on lean startup and the scientific approach to
entrepreneurship. Note that the strongest study measures better \emph{stopping}
decisions rather than higher returns --- which is a real benefit, and not the
one usually claimed.}
\label{tab:startup-lean-evidence}
\begin{tblr}{
  width=\textwidth,
  colspec={Q[l,0.21\textwidth] X[l] Q[l,0.17\textwidth]},
  row{1}={font=\bfseries,bg=AlmanachMist},
  hlines,vlines}
Study & Finding & Design \\
Camuffo et al.\ (2020)\textsuperscript{\cite{camuffoScientificApproachEntrepreneurial2020}}
  & Founders trained in hypothesis-driven decision making were more likely to
  \emph{terminate} projects and to \emph{pivot radically}
  & Randomised controlled trial, 116 Italian startups \\
Camuffo et al.\ (2024)\textsuperscript{\cite{camuffoScientificApproachReplication2024}}
  & Replication confirms the effect on termination; evidence of a
  \emph{non-linear} effect on pivots --- a few beat none or many
  & Large-scale replication and extension \\
Welter et al.\ (2021)\textsuperscript{\cite{welterRoadEntrepreneurialSuccess2021}}
  & Collecting preorders ($\beta=0.89$, $p=0.04$ for growth) and pivoting on
  feedback ($\beta=0.34$, $p=0.03$) helped. \textbf{Prototyping and
  experimentation showed no significant effect} --- and writing a business plan
  \emph{did} ($\beta=0.30$, $p=0.036$)
  & Survey, 120 US entrepreneurs \\
Felin et al.\ (2020)\textsuperscript{\cite{felinLeanStartupBusiness2020}}
  & Theoretical critique: the method says little about where hypotheses
  \emph{come from}, and customer feedback on a nascent product may be a weak
  signal that biases toward incremental outcomes
  & Conceptual \\
\end{tblr}
\end{table}

\begin{engineernote}[Reading the evidence honestly]
Three conclusions, none of them the marketing version.

\emph{1. The measured benefit is better decisions, especially about stopping.}
The randomised trial's clearest effect is that trained founders killed bad
projects sooner. That is genuinely valuable --- an entrepreneur's scarcest
asset is the time they are spending on the wrong thing --- but it is not the
same as ``lean startup makes startups succeed''.

\emph{2. The parts are not equally supported.} Welter and colleagues found
preorders and evidence-driven pivots associated with growth while prototyping
and experimentation were not. Treat the method as a menu of practices with
different evidential standing, not as a doctrine to adopt whole.

\emph{3. The plan-versus-lean dichotomy is false.} In the same study, writing
a business plan was \emph{also} associated with success. The popular reading
that plans are waste is not supported; what is supported is that a plan
containing untested assumptions should be treated as a list of experiments
rather than a commitment.
\end{engineernote}

\begin{cautionbox}[The critique worth taking seriously]
Felin and colleagues make an argument that no amount of iteration
answers\textsuperscript{\cite{felinLeanStartupBusiness2020}}: build--measure--learn
tells you how to \emph{test} a hypothesis and almost nothing about where a
hypothesis worth testing comes from. Experimentation is a search procedure,
and a search procedure guided by customer feedback on an early product hill-
climbs toward local improvements. It is well suited to finding the next
increment and poorly suited to finding a non-obvious position that customers
could not have described.

The practical reading is not to abandon the method but to recognise its
division of labour: \emph{theory generates the hypothesis; experiment
adjudicates it}. Ries's loop is excellent at the second job. The first still
requires a point of view about the world that no customer interview will hand
you --- which is precisely why the value-hypothesis question in this chapter's
opening section comes before any experiment is designed.
\end{cautionbox}

\subsection{Adapting the Method to AI Products}

Lean startup was written for software whose behaviour is deterministic and
whose cost falls to nearly zero at scale. Neither holds for AI products, and
four adjustments follow.

\begin{table}[htbp]
\centering
\caption{Where the standard method needs modification for an AI product, and
what to do instead. The pattern: AI moves cost and risk from the build phase
into the measure phase.}
\label{tab:startup-lean-for-ai}
\begin{tblr}{
  width=\textwidth,
  colspec={Q[l,0.20\textwidth] X[l] X[l]},
  row{1}={font=\bfseries,bg=AlmanachMist},
  hlines,vlines}
Assumption & Why it breaks for AI & Adaptation \\
The MVP is cheap to build & Model quality is not a feature you can stub; it
  is the product & Use a Wizard-of-Oz or concierge MVP --- humans deliver the
  output --- to test demand before any model exists \\
Output quality is binary & Quality is a distribution, and the tail is what
  users remember & Declare the acceptance threshold and the failure behaviour
  as part of the hypothesis \\
Marginal cost is near zero & Inference costs money per request, every request
  & Test the unit economics as an explicit hypothesis
  (\cref{sec:startup-unit-economics}), not as a later optimisation \\
Faster iteration is always better & Model changes shift behaviour globally,
  not locally & Keep a frozen evaluation set; a change that improves the demo
  may regress cases you already sold \\
\end{tblr}
\end{table}

\begin{engineernote}
The Wizard-of-Oz MVP deserves emphasis, because it is unusually well suited to
AI and unusually underused. If you are proposing to spend six months building
a model to draft responses, first have two people draft those responses by
hand behind the same interface for thirty customers. You will learn whether
anyone wants the output, what ``good'' means to them, and what the acceptance
threshold is --- and the transcripts become your evaluation set and your
training data. The manual version costs weeks rather than months, and it
answers the question the model cannot: whether the problem was worth solving.
\end{engineernote}

\begin{practicallab}[Run one turn of the loop, backwards]
\textbf{Goal.} Convert a product idea into a completed experiment, and observe
how much you learn without building the thing.

\textbf{Protocol.}
\begin{enumerate}
  \item Write your value hypothesis and growth hypothesis as separate,
        falsifiable sentences. Each must name a number and a threshold.
  \item Rank your assumptions by \emph{leap of faith} --- how much of the
        business dies if this one is false --- not by how easy they are to test.
  \item Take the riskiest. Design the loop \emph{backwards}: what must you
        learn, what would you measure, what is the smallest instrument that
        produces the measurement?
  \item Choose the MVP form: smoke test, concierge, or Wizard of Oz. Justify
        why the cheaper option would not answer the question.
  \item \textbf{Pre-declare the decision rule}: at what result do you persevere,
        pivot, or stop? Write it down before running.
  \item Run it with at least $20$ real people who do not know you.
  \item Report by cohort, never as a cumulative total.
  \item Hold the pivot-or-persevere meeting and apply the rule from step~5
        exactly as written.
\end{enumerate}

\textbf{Deliverable.} The two hypotheses, the ranked assumption list, the
pre-declared decision rule, the cohort table, and the decision taken.

\textbf{The point.} Step~5 is the phase. A decision rule written after seeing
the data is not a decision rule, and this is where most teams quietly
substitute the outcome they wanted.
\end{practicallab}

\begin{checkpoint}
Your AI assistant has $12{,}000$ signups, up from $2{,}000$ last quarter, and
the team wants to raise on that number. Name the metric that would falsify the
optimistic reading, state which of the two hypotheses each of the two numbers
speaks to, and give the pivot type you would consider first if per-cohort
retention is falling.
\end{checkpoint}

\section{The Evidence Loop}

\begin{figure}[htbp]
\centering
\begin{tikzpicture}[
    node distance=7mm and 9mm,
    venture step/.style={draw,rounded corners=1mm,fill=LimeGreen!8,
        minimum width=2.7cm,minimum height=8mm,align=center,
        font=\sffamily\scriptsize},
    venture arrow/.style={-{Stealth[length=2mm]},thick}]
\node[venture step] (assumption) {Rank the riskiest\\assumption};
\node[venture step,right=of assumption] (test)
    {Design the smallest\\credible test};
\node[venture step,below=of test] (measure)
    {Measure behaviour,\\quality, cost, and risk};
\node[venture step,left=of measure] (decide) {Continue, change,\\or stop};
\draw[venture arrow] (assumption) -- (test);
\draw[venture arrow] (test) -- (measure);
\draw[venture arrow] (measure) -- (decide);
\draw[venture arrow] (decide) -- (assumption);
\end{tikzpicture}
\caption{The start-up evidence loop turns assumptions into explicit tests and
decision gates.}
\label{fig:startup-evidence-loop}
\end{figure}

\begin{v3topic}{Assumption Ledger}
Maintain one row per important assumption.
Record the consequence of being wrong, present evidence, test method, decision
threshold, owner, and review date.
Rank by \emph{uncertainty multiplied by consequence}, not by which experiment
is most entertaining to build.

\end{v3topic}
\begin{v3topic}{Experiment Before Architecture}
Match the test to the hypothesis.
Use interviews and observation for workflow pain, a concierge service for
value, a frozen representative dataset for technical feasibility, a paid
pilot for willingness to pay, and a threat or compliance review for risks
that cannot ethically be discovered in production.

\end{v3topic}
\begin{v3topic}{Evidence Gates}
Stage gates are useful when they specify evidence rather than ceremony.
Cooper emphasises adaptable, accountable gates with explicit criteria instead
of a rigid linear process\textsuperscript{\cite{cooperPerspectiveStageGateIdeatoLaunch2008}}.
For an AI system, every gate should include product, model, operations, safety,
security, privacy, and economic evidence.
\visualseparator
\begin{tblr}{
    width=\linewidth,
    colspec={Q[l,0.16\linewidth]X[l]X[l]X[l]},
    row{1}={font=\bfseries,bg=LimeGreen!15},
    hlines,
    vlines}
Gate & Question & Evidence & Stop or change when \\
Discovery & Is the problem frequent and costly? & Observed workflows,
interviews, current workaround & No repeated behaviour or accountable buyer \\
Feasibility & Can a simple system meet the target? & Baseline, representative
split, error slices, cost profile & Target requires unavailable data or unsafe
behaviour \\
Pilot & Does the workflow improve in context? & Prospective usage, overrides,
time saved, incident log & Users route around the system or errors concentrate \\
Production & Can the service be operated responsibly? & SLOs, monitoring,
security review, legal basis, fallback & Residual risk or unit economics exceed
tolerance \\
Scale & Does value persist across customers? & Cohort retention, support load,
shift tests, margin by segment & Custom work grows faster than repeatable value \\
\end{tblr}
\captionof{table}{Evidence gates for an AI product from discovery to scale.}
\label{tab:startup-evidence-gates}
\end{v3topic}

\section{AI Feasibility and Product Architecture}

\begin{v3topic}{Define the Product Contract}
Specify the input, output, user, decision, latency, availability, acceptable
error, abstention path, privacy boundary, and operating conditions.
Then write the evaluation contract: dataset provenance, split logic, metrics,
slices, uncertainty, baselines, and release threshold.
The product and evaluation contracts must describe the same real-world use.

\end{v3topic}
\begin{v3topic}{Build, Buy, or Partner}
\textbf{Build} when proprietary capability is central and evidence supports
the long-term cost.
\textbf{Buy} when a service is substitutable and contracts meet privacy,
security, availability, and exit needs.
\textbf{Partner} when success requires domain access or distribution that the
team cannot recreate.
Preserve an exit path: exportable data, portable evaluations, versioned
prompts or adapters, and a non-proprietary fallback.

\end{v3topic}
\begin{v3topic}{Where Defensibility Comes From}
Model access alone is rarely a durable moat.
More defensible assets are permissioned data generated by a valuable workflow,
hard-won integration, trust and distribution, a fast learning loop, and
evidence that the system performs safely in a demanding niche.
Feedback data is useful only when its labels are meaningful and its collection
does not create harmful incentives.
\end{v3topic}

\begin{engineernote}[Prototype the uncertainty boundary]
A polished interface can hide unreliable behaviour.
During discovery, show confidence limits, abstentions, citations, provenance,
and failure cases explicitly.
The team needs to learn where automation stops being dependable before users
learn that lesson in production.
\end{engineernote}

\section{Unit Economics for AI Systems}
\label{sec:startup-unit-economics}

For one completed customer workflow, a useful contribution estimate is
\begin{equation}
  M = P - C_{\mathrm{inference}} - C_{\mathrm{data}}
      - C_{\mathrm{review}} - C_{\mathrm{support}} - C_{\mathrm{risk}},
\end{equation}
where $P$ is realised revenue and the costs include failed requests and human
review, not merely successful model calls.
If a workflow needs $n$ attempts with per-attempt cost $c$ and succeeds with
probability $p$, the expected inference cost per successful workflow is
approximately
\begin{equation}
  C_{\mathrm{inference/success}} = \frac{n c}{p}.
\end{equation}
This approximation must be expanded when retries are correlated, cache hits
matter, long contexts change cost, or expensive failures create support and
liability exposure.

\begin{workedexample}[The cheaper model is not always the cheaper product]
Model A costs \$0.04 per attempt and completes 80\% of workflows without human
review.
Model B costs \$0.10 and completes 95\%.
A failed workflow costs \$1.20 in review time.
Ignoring other costs, expected cost is
$0.04 + 0.20(1.20)=\$0.28$ for A and
$0.10 + 0.05(1.20)=\$0.16$ for B.
Optimise the full workflow, then verify the assumptions with measured queues
and failure categories.
\end{workedexample}

\section{Responsible Commercialisation}

\begin{v3topic}{Data and Rights}
Record the origin, licence, consent or legal basis, purpose, retention, and
deletion route for each data source.
Separate the right to access data from the right to train a model, generate an
output, or transfer data to a supplier.

\end{v3topic}
\begin{v3topic}{Claims and Sales}
Marketing claims are engineering requirements in public clothing.
Attach scope, date, dataset, uncertainty, and known limits to every material
performance claim.
Do not convert a laboratory result into claims about safety, health, finance,
or legal compliance that the study did not test.

\end{v3topic}
\begin{v3topic}{Operating Responsibility}
Name owners for model changes, security patches, incidents, customer notices,
data-subject requests, and rollback.
Budget for monitoring, red-team work, documentation, insurance, and specialist
review where the domain requires it.

\end{v3topic}
\begin{v3topic}{Kill Criteria Are an Asset}
Stopping preserves time and credibility.
Precommit to kill or pivot criteria for absent demand, unattainable quality,
unsafe residual risk, unlawful data dependence, integration burden, and unit
economics that do not improve with realistic scale.
\end{v3topic}

\section{Practical Lab: One-Page AI Product Brief}

\begin{practicallab}[Turn a model idea into an evidence plan]
Choose one workflow and produce a one-page brief containing:
\begin{enumerate}
    \item the user, triggering event, current workaround, and measurable pain;
    \item the non-AI baseline and the smallest AI-assisted intervention;
    \item product and evaluation contracts, including error costs and
          abstention;
    \item the five riskiest assumptions and one cheap, ethical test for each;
    \item a cost model per successful workflow and a sensitivity analysis;
    \item data rights, safety, security, privacy, and legal review questions;
    \item pilot, production, and scale gates with owners and kill criteria.
\end{enumerate}
Interview three relevant users and revise the brief using observed behaviour.
The deliverable is the changed assumption ledger, not a slide deck claiming
that every interview was encouraging.
\end{practicallab}

\begin{checkpoint}
\begin{enumerate}
    \item Which of the four hypotheses does a high offline score support?
    \item What is the strongest non-AI baseline for your proposed workflow?
    \item Which cost grows when model quality falls even if token cost stays
          fixed?
    \item What evidence would make you stop the project next week?
\end{enumerate}
\end{checkpoint}

%
\section{Deep Dive: Portfolio Evidence and Option Value}
\begin{v3topic}{Fund Learning Milestones}
Tie spending to the next uncertainty-reducing milestone: verified problem,
representative data, baseline feasibility, paid workflow adoption, responsible
operation, or repeatable economics.
Do not confuse accumulated code with reduced venture risk.

\end{v3topic}
\begin{v3topic}{Preserve Strategic Options}
Design pilots so data, evaluations, interfaces, and customer learning remain
useful if the model, supplier, or route to market changes.
Portability and a credible manual fallback create option value when technology
and regulation move faster than the plan.
\end{v3topic}

\chapterreview{Problem}{Experiment}{Economics}{Scale}

\section{Further Reading}

\begin{v3topic}{Lean startup: the primary text and the evidence}
\begin{itemize}
  \item \fullcite{riesLeanStartup2011} --- the founding text; the German
    edition is \fullcite{riesLeanStartupDeutsch2014}.
  \item \fullcite{blankWhyLeanStartup2013} --- the customer-development
    tradition lean startup grew out of, in twelve pages.
  \item \fullcite{osterwalderBusinessModelGeneration2010} --- the business
    model canvas, the usual companion tool for stating hypotheses.
  \item \fullcite{camuffoScientificApproachEntrepreneurial2020} --- the
    randomised controlled trial; read this before quoting any claim about
    whether the method works.
  \item \fullcite{welterRoadEntrepreneurialSuccess2021} --- which specific
    practices correlate with success, and which do not.
  \item \fullcite{felinLeanStartupBusiness2020} --- the substantive critique:
    where do hypotheses come from?
\end{itemize}
\end{v3topic}

Customer interviews and evidence-based discovery are developed in
\cite{constableTalkingHumansSuccess2014}.
Cooper's stage-gate work explains how explicit, adaptable gates support
accountability without forcing innovation into a rigid waterfall
\cite{cooperPerspectiveStageGateIdeatoLaunch2008}.
